% Part: normal-modal-logic
% Chapter: syntax-and-semantics
% Section: normal-modal-logics

\documentclass[../../../include/open-logic-section]{subfiles}

\begin{document}

\olfileid{nml}{syn}{nor}

\olsection{স্বাভাবিক মোডাল যুক্তিবিদ্যা}

একটি মোডাল যুক্তিবিদ্যাকে উপযুক্তভাবে সুবিন্যস্ত !!{formula}s-এর যেকোনো
সমষ্টি হিসেবে দেখা যায়। কিছু মোডাল যুক্তিবিদ্যা বিশেষ আগ্রহের: তাদের
আকর্ষণীয় ব্যাখ্যা ও প্রয়োগ আছে, তবে প্রধান কারণ হলো সেগুলি তাত্ত্বিক
দিক থেকে ভালোভাবে বোঝা যায়। এগুলিই তথাকথিত স্বাভাবিক মোডাল
যুক্তিবিদ্যা।

\begin{defn}\ollabel{defn:modal-logic}
একটি \emph{মোডাল যুক্তিবিদ্যা}~$\Log L$ হলো মৌলিক মোডাল ভাষার
!!{formula}s-এর এমন একটি সমষ্টি, যা
\begin{enumerate}
\item সমস্ত প্রতিজ্ঞামূলক সর্বতঃসত্য সূত্র ধারণ করে,
\item অভিন্ন প্রতিস্থাপনের অধীনে বদ্ধ থাকে, এবং
\item মোডাস পোনেন্সের অধীনে বদ্ধ থাকে, অর্থাৎ $!A$ এবং $(!A \lif
  !B) \in \Log L$ হলে $!B \in \Log L$-ও হয়।
\end{enumerate}
\end{defn}

\begin{ex}
নিম্নলিখিত !!{formula}s-এর সমষ্টিগুলি খুব আকর্ষণীয় না হলেও এই
সংজ্ঞা মেনে চলে এবং তাই মোডাল যুক্তিবিদ্যা হিসেবে গণ্য হয়:
\begin{enumerate}
\item সমস্ত সর্বতঃসত্য সূত্রের সমষ্টি
\item সমস্ত !!{formula}s-এর সমষ্টি (অসঙ্গত মোডাল যুক্তিবিদ্যা)
\end{enumerate}
\end{ex}

\begin{defn}\label{defn:normal-modal-logic}
একটি মোডাল যুক্তিবিদ্যা~$\Log L$ \emph{স্বাভাবিক} তখনই এবং কেবল
তখনই, যখন
\begin{enumerate}
\item এতে নিচের দুটি সূত্রই থাকে:
\begin{align}
\tag{K} \Box(p \lif q) & \lif (\Box p \lif \Box q) \\
\tag{Dual} \Diamond p & \liff \lnot \Box \lnot p
\end{align}
\item এবং এটি আবশ্যকীকরণের অধীনে বদ্ধ থাকে, অর্থাৎ $!A \in
  \Log L$ হলে $\Box !A \in \Log L$-ও হয়।
\end{enumerate}
\end{defn}

মোডাল যুক্তিবিদ্যা সংজ্ঞায়িত করার দুটি প্রধান পদ্ধতি আমরা দেখব।
প্রথমটি প্রমাণতাত্ত্বিক: নির্দিষ্ট !!{derivation} পদ্ধতিতে
নিষ্পাদনযোগ্য !!{formula}s-এর সমষ্টি হিসেবে। দ্বিতীয়টি মডেলতাত্ত্বিক:
নির্দিষ্ট মডেল-শ্রেণিতে বৈধ !!{formula}s-এর সমষ্টি হিসেবে। সাধারণ
প্রতিজ্ঞামূলক যুক্তিবিদ্যায় (কিংবা প্রথম-ক্রম যুক্তিবিদ্যাতেও)
পরিস্থিতি একই রকম। সেখানেও !!{formula}s-এর একটি সমষ্টিকে দুইভাবে
চিহ্নিত করা যায়: যেমন সব সত্যমান-আরোপে সত্য সর্বতঃসত্য সূত্রের
সমষ্টি হিসেবে এবং কোনো !!{derivation} পদ্ধতিতে সমস্ত !!{derivable}
!!{formula}s-এর সমষ্টি হিসেবে। স্বাভাবিক মোডাল যুক্তিবিদ্যার ক্ষেত্রে
সম্পর্কমূলক মডেল ও স্বতঃসিদ্ধমূলক (হিলবার্ট-ধরনের) প্রমাণপদ্ধতি
ব্যবহার করে এই দুইভাবে যুক্তিবিদ্যা নির্দিষ্ট করা যায়।

\end{document}
